\section*{Прикладная задача}
Дано $N$ независимых работ, для каждой работы задано время выполнения.
Требуется построить расписание выполнения работ без прерываний на $M$
процессорах. На расписании должно достигаться минимальное значение
\textit{критерия К2}.

\textbf{Критерий К2:} суммарное время ожидания (т.е. сумма, по всем работам в
расписании, времён завершения работ)
\section*{Формальная постановка задачи}
\textbf{Дано:}

\begin{itemize}
  \item $N$ – количество работ;
  \item $J = \{ j_1, j_2, \ldots, j_N \}$ – множество работ;
  \item $\tau = \{ t_1, t_2, \ldots, t_N \}$ – множество времён выполнения соответствующих заданий $j_i$, 
  где $\forall i \in [1, N]: t_i > 0$;
  \item $M$ – количество процессоров;
  \item $P = \{ p_1, p_2, \ldots, p_M \}$ – множество процессоров, на которых выполняются работы.
\end{itemize}

\textbf{Расписание}

В данной работе расписание представляется набором
упорядоченных списков работ, по одному для каждого процессора:

\[
G = \{ G_1, G_2, \ldots, G_M \},
\]

где каждый \(G_j\) — последовательность индексов работ, выполняющихся
на процессоре \(p_j\) в порядке их выполнения:

\[
G_j = [\, i_1,\, i_2,\, \ldots,\, i_{k_j} \,].
\]

Каждая работа встречается ровно в одном списке:
\[
\bigcup_{j=1}^{M} G_j = \{1,2,\ldots,N\}, \qquad
G_j \cap G_k = \varnothing \text{ при } j \neq k.
\]

Для совместимости может использоваться бинарная матрица
\(H \in \{0,1\}^{N\times M}\), определяемая как
\(h_{ij}=1\) тогда и только тогда, когда \(i \in G_j\).
Матрица \(H\) вспомогательна и не задаёт порядок выполнения работ.


\textbf{Требуется}

Необходимо построить расписание \(G = \{G_1,\ldots,G_M\}\), минимизирующее
суммарное время завершения всех работ (критерий \(K_2\)).


\textbf{Критерий минимизации:}

Выбран \textbf{критерий $K_2$ — суммарное время ожидания.}
  
Для каждой работы $i \in G_j$ определим время завершения $C_i$ как сумму длительностей всех работ,
выполняющихся на данном процессоре до и включая текущую:

\[
C_i = \sum_{k \in G_j, \; k \text{ перед } i} t_k + t_i
\]

Тогда суммарное время ожидания всех работ (критерий $K_2$) определяется как:

\[
K_2 = \sum_{i=1}^{N} C_i
\]

Необходимо найти такое распределение работ по процессорам (расписание $HP$),
которое минимизирует данный критерий:

\[
\min_{HP} K_2 = \min_{HP} \sum_{i=1}^{N} C_i
\]

\textbf{Ограничения:}

\begin{itemize}
  \item Каждый процессор $p_j \in P$ в любой момент времени может выполнять не более одного задания.
  \item Во время выполнения задания прерывания не допускаются.
  \item Процессор может мгновенно переключаться между заданиями без потерь времени.
  \item Время выполнения каждой работы $t_i \in \tau$ фиксировано и известно заранее.
\end{itemize}